\documentclass{ctexart}
\usepackage{template}
\usepackage{extarrows}

\begin{document}\pagestyle{empty}
\begin{center}\Large
    北京大学数学科学学院2024-25高等数学B1期中模拟
\end{center}
\begin{problem}[1.(15\songti{分})]
    解答下列问题.
    \begin{enumerate}[label=\textbf{(\arabic*)}]
        \item \textbf{(5\songti{分})}\ 求序列极限$$\lim_{n\to\infty}\left|\cos\left(\sqrt{n^2+1}\pi\right)\right|$$
        \item \textbf{(5\songti{分})}\ 求函数极限$$\lim_{x\to+\infty}\left(\dfrac{x^2-1}{x^2+1}\right)^{x^2}$$
        \item \textbf{(5\songti{分})}\ 求函数$$f(x)=\left\{\begin{array}{l}x^2\sin\dfrac{1}{x},x\neq 0\\0,x=0\end{array}\right.$$的导函数$f'(x)$.
    \end{enumerate}
\end{problem}
\begin{solution}
    \begin{enumerate}[label=\textbf{(\arabic*)}]
        \item $$\begin{aligned}
                \lim_{n\to\infty}\left|\cos\left(\pi\sqrt{n^2+1}\right)\right|
                &= \lim_{n\to\infty}\left|\cos\left(\pi\sqrt{n^2+1}-n\pi\right)\right| \\
                &= \lim_{n\to\infty}\left|\cos\dfrac{\pi}{\sqrt{n^2+1}+n}\right| \\
                &= \cos 0 \\
                &= 1
            \end{aligned}$$
        \item $$\begin{aligned}
                \lim_{x\to+\infty}\left(\dfrac{x^2-1}{x^2+1}\right)^{x^2}
                &= \lim_{x\to\infty}\left(1-\dfrac{2}{x^2+1}\right)^{\frac{x^2+1}{2}\cdot\frac{2x^2}{x^2+1}} \\
                &= \e^{-2}
            \end{aligned}$$
        \item 当$x\neq0$时有
            $$\dfrac{\di f(x)}{\dx}=2x\sin\dfrac{1}{x}+x^2\left(-\dfrac{1}{x^2}\cos\dfrac{1}{x}\right)=2x\sin\dfrac{1}{x}-\cos\dfrac{1}{x}$$
            又$$f'(0)=\lim_{x\to0}\dfrac{f(x)-f(0)}{x}=\lim_{x\to0}\dfrac{x^2\sin\dfrac{1}{x}}{x}=\lim_{x\to0}x\sin\dfrac{1}{x}=0$$
            于是$$f'(x)=\left\{\begin{array}{l}
                2x\sin\dfrac{1}{x}-\cos\dfrac{1}{x},x\neq0\\0,x=0
            \end{array}\right.$$
    \end{enumerate}
\end{solution}
\begin{problem}[2.(10\songti{分})]
    设参数方程$$\left\{\begin{array}{l}x=\e^t\sin2t\\y=\e^t\cos t\end{array}\right.$$确定函数$y=f(x)$.求$\dfrac{\di y}{\dx},\dfrac{\di^2y}{\dx^2}$.
\end{problem}
\begin{solution}
记$D(t)=\sin2t+2\cos2t$.由参数方程求导,在$D(t)\neq0$的区间上有
    \[\dfrac{\di y}{\dx}=\dfrac{\cos t-\sin t}{D(t)}.\]
    再次求导得到
    \[\dfrac{\di^2y}{\dx^2}
    =\dfrac{-(\sin t+\cos t)D(t)-(\cos t-\sin t)(2\cos2t-4\sin2t)}{\e^tD(t)^3}.\]
\end{solution}
\begin{problem}[3.(10\songti{分})]
    求不定积分$$\int\dfrac{x^3+1}{x(x-1)^3}\dx$$
\end{problem}
\begin{solution}
    置$u=x-1$,于是
    $$\begin{aligned}
        \int\dfrac{x^3+1}{x(x-1)^3}\dx
        &= \int\dfrac{u^3+3u^2+3u+2}{u^3(u+1)}\di u
    \end{aligned}$$
    设$$\dfrac{u^3+3u^2+3u+2}{u^3(u+1)}=\dfrac{A}{u+1}+\dfrac{B}{u}+\dfrac{C}{u^2}+\dfrac{D}{u^3}$$
    于是$$\left\{\begin{array}{l}
        A+B=1\\
        B+C=3\\
        C+D=3\\
        D=2
    \end{array}\right.$$
    解得$A=-1,B=2,C=1,D=2$.
    于是$$\begin{aligned}
        \int\dfrac{x^3+1}{x(x-1)^3}\dx
        &= \int\left(\dfrac{2}{u^3}+\dfrac{1}{u^2}+\dfrac{2}{u}-\dfrac{1}{u+1}\right)\di u\\
        &= -\dfrac{1}{(x-1)^2}-\dfrac{1}{x-1}+2\ln\left|x-1\right|-\ln\left|x\right|+C
    \end{aligned}$$
\end{solution}
\begin{problem}[4.(10\songti{分})]
    求心形线$$r=a(1+\cos\theta),a>0,0\leqslant\theta\leqslant2\pi$$的弧长.
\end{problem}
\begin{solution}
根据极坐标弧长公式,
    \[\begin{aligned}
      L&=\int_0^{2\pi}\sqrt{r^2+(r')^2}\di\theta
      =2a\int_0^{2\pi}|\cos(\theta/2)|\di\theta\\
       &=4a\int_0^\pi|\cos u|\di u=8a.
    \end{aligned}\]
\end{solution}
\begin{problem}[5.(10\songti{分})]
    设$t\in(0,1)$,$f(t)$表示曲线$y=\sec x$与直线$x=0,y=0$和$x=\arcsin t$围成的封闭图形的面积.
    \begin{enumerate}[label=\textbf{(\arabic*)}]
        \item \textbf{(5\songti{分})}\ 求$f(t)$的导数$f'(t)$.
        \item \textbf{(5\songti{分})}\ 求$f(t)$.
    \end{enumerate}
\end{problem}
\begin{solution}
    \begin{enumerate}[label=\textbf{(\arabic*)}]
        \item 由题意可得$$f(t)=\int_{0}^{x(t)}y\dx=\int_0^{\arcsin{t}}\sec x\dx$$
            于是$$f'(t)=\sec(\arcsin t)\cdot\left(\arcsin t\right)'=\dfrac{1}{\sqrt{1-t^2}\sqrt{1-t^2}}=\dfrac{1}{1-t^2}$$
        \item 由\textbf{(1)}的结果可得
            $$\int f'(t)\di t=\int\dfrac{\di t}{1-t^2}=\dfrac{1}{2}\int\left(\dfrac{1}{t+1}-\dfrac{1}{t-1}\right)\di t=\dfrac{1}{2}\ln\left|\dfrac{t+1}{t-1}\right|+C$$
            将$f$连续延拓至$t=0$,则$$f(0)=\int_0^{\arcsin0}\sec x\dx=0$$
            于是$$f(t)=\dfrac{1}{2}\ln\left|\dfrac{t+1}{t-1}\right|$$
    \end{enumerate}
\end{solution}
\begin{problem}[6.(10\songti{分})]
    设函数$f(x)$在$[0,1]$上连续,且有$$\int_0^1f(x)\dx=\int_0^1xf(x)\dx=0$$试证明:$f(x)$在$[0,1]$上至少有两个零点.
\end{problem}
\begin{solution}
若$f\equiv0$,结论显然成立.否则,由$\int_0^1f=0$及连续性,$f$必须同时取正值和负值,从而在$(0,1)$内有零点.
    假设恰有一个零点$c\in(0,1)$.由连续性,$f$在$(0,c)$和$(c,1)$上分别保持固定符号,且两侧符号相反,否则$\int_0^1f\neq0$.
    因此$(x-c)f(x)$在除$c$外的整个区间上保持同一个严格符号,所以
    \[\int_0^1(x-c)f(x)\dx\neq0.\]
    但题设又给出该积分等于$\int_0^1xf(x)\dx-c\int_0^1f(x)\dx=0$,矛盾.故至少有两个不同零点.
\end{solution}
\begin{problem}[7.(20\songti{分})]
    解答下列问题.
    \begin{enumerate}[label=\textbf{(\arabic*)}]
        \item \textbf{(8\songti{分})}\ 设序列$\left\{x_n\right\}$满足$\displaystyle x_n=\sum_{i=1}^n\dfrac{1}{i}-\ln n$,试证明$\displaystyle\lim_{n\to\infty}x_n$存在且有限.
        \item \textbf{(12\songti{分})}\ 对于$n\in\N$,定义$$I_n=\int_0^\frac{\pi}{2}\dfrac{\sin^2nt}{\sin t}\di t$$求序列极限$\displaystyle\lim_{n\to\infty}\dfrac{I_n}{\ln n}$.
    \end{enumerate}
\end{problem}
\begin{solution}
    \begin{enumerate}[label=\textbf{(\arabic*)}]
        \item 置$f(x)=\dfrac{1}{x}$,于是$$\forall k\in\N^*,\forall x\in(k,k+1),\dfrac{1}{k}>f(x)>\dfrac{1}{k+1}$$
            对上式积分有$$\dfrac{1}{k+1}<\int_{k}^{k+1}f(x)\dx<\dfrac{1}{k}$$
            对上式求和有$$\sum_{k=1}^{n-1}\dfrac{1}{k+1}<\int_1^nf(x)\dx<\sum_{k=1}^{n-1}\dfrac{1}{k}$$
            于是$$x_n+\ln n-1<\left.(\ln x)\right|_1^n<x_n+\ln n-\dfrac{1}{n}$$
            所以$$\dfrac{1}{n}<x_n<1\quad(n\geqslant2)$$
            从而$\left\{x_n\right\}$有界.我们又有
            $$x_{n+1}-x_n=\dfrac{1}{n+1}+\ln\dfrac{n}{n+1}=\ln\dfrac{n}{n+1}+1-\dfrac{n}{n+1}<0$$
            于是$\left\{x_n\right\}$递减,从而$\displaystyle\lim_{n\to\infty}x_n$存在.
        \item 不难发现$$I_1=\int_0^{\frac{\pi}{2}}\sin t\di t=1$$
            而
            $$\begin{aligned}
                I_n-I_{n-1}
                &= \int_0^{\frac{\pi}{2}}\dfrac{\sin^2nt-\sin^2(n-1)t}{\sin t}\di t \\
                &= \int_0^{\frac{\pi}{2}}\dfrac{\sin(n+(n-1))t\cdot\sin(n-(n-1))t}{\sin t}\di t \\
                &= \int_0^{\frac{\pi}{2}}\sin(2n-1)t\di t \\
                &= \dfrac{1}{2n-1}\int_0^{n\pi-\frac{\pi}{2}}\sin u\di u \\
                &= \dfrac{1}{2n-1}
            \end{aligned}$$
            于是$$I_n=1+\dfrac{1}{3}+\dfrac{1}{5}+\cdots+\dfrac{1}{2n-1}$$
            设$\displaystyle t_n=x_n+\ln n=\sum_{i=1}^n\dfrac{1}{i}$,于是$I_n=t_{2n}-\dfrac{1}{2}t_n$.
            于是$$\begin{aligned}
                \dfrac{I_n}{\ln n}
                &= \dfrac{x_{2n}+\ln2n}{\ln n}-\dfrac{x_n+\ln n}{2\ln n} \\
                &= \dfrac{x_{2n}}{\ln n}-\dfrac{x_n}{2\ln n}+\dfrac{\ln 2}{\ln n}+\dfrac{1}{2}
            \end{aligned}$$
            注意到$\forall n\in\N^*,0<x_n\leqslant1$.于是$$\lim_{n\to\infty}\dfrac{I_n}{\ln n}=\dfrac{1}{2}$$
    \end{enumerate}
\end{solution}
\begin{problem}[8.(15\songti{分})]
    设函数$f(x)$在$[0,\pi]$上连续.对于$n\in\N$,试证明$$\lim_{n\to\infty}\int_{0}^{\pi}f(x)\left|\sin(nx)\right|\dx=\dfrac{2}{\pi}\int_0^\pi f(x)\dx$$
\end{problem}
\begin{solution}
将$[0,\pi]$分成$n$个区间$[k\pi/n,(k+1)\pi/n]$,其中$k=0,\ldots,n-1$.
    由积分第一中值定理,每段存在$\xi_k$属于该区间,使
    \[\int_{k\pi/n}^{(k+1)\pi/n}f(x)|\sin(nx)|\dx
      =f(\xi_k)\int_{k\pi/n}^{(k+1)\pi/n}|\sin(nx)|\dx=\dfrac2n f(\xi_k).\]
    因而由Riemann和的定义,
    \[\begin{aligned}
      \lim_{n\to\infty}\int_0^\pi f(x)|\sin(nx)|\dx
      &=\dfrac2\pi\lim_{n\to\infty}\sum_{k=0}^{n-1}f(\xi_k)\dfrac\pi n\\
      &=\dfrac2\pi\int_0^\pi f(x)\dx.
    \end{aligned}\]
\end{solution}
\end{document}
